\documentclass{article}
\usepackage[utf8]{inputenc}
\usepackage{amsmath}
\usepackage{geometry}
\geometry{margin=2cm}
\begin{document}

\section*{Análise da Questão - ENEM 2024 - Questão 175}

\subsection*{Contexto}
Uma indústria utiliza folhas retangulares de alumínio de dimensões 10 cm por 20 cm para fabricar embalagens cilíndricas. Utiliza-se a folha para formar a superfície lateral do cilindro, fechando-se com tampo e fundo circulares. As duas opções surgem dependendo de qual lado da folha è usado como altura do cilindro e qual è a circunferência da base.

\subsection*{Solução}
Para determinar a maior capacidade, calculamos os volumes dos dois cilindros possíveis:

\begin{itemize}
  \item \textbf{Embalagem 1:}
    \begin{itemize}
      \item Altura $h = 10$ cm
      \item Circunferência da base $C=20$ cm
      \item Raio $r = \frac{20}{2\pi} = \frac{10}{\pi}$
      \item Volume $V = \pi r^2 h = \pi \left(\frac{10}{\pi}\right)^2 10 = 1000 \pi$
    \end{itemize}
  \item \textbf{Embalagem 2:}
    \begin{itemize}
      \item Altura $h = 20$ cm
      \item Circunferência da base $C=10$ cm
      \item Raio $r = \frac{10}{2\pi} = \frac{5}{\pi}$
      \item Volume $V = \pi r^2 h = \pi \left(\frac{5}{\pi}\right)^2 20 = 500 \pi$
    \end{itemize}
\end{itemize}

\no volume maior è $1000\pi$ cm$^3$, referente à embalagem 1.

\subsection*{Conclusão}
O volume da embalagem com maior capacidade è:

\[
\boxed{1000 \pi \text{ cm}^3}.
\]

Esta resposta está de acordo com o gabarito (alternativa D).

\end{document}